\documentclass{article}

\usepackage{fancyhdr}
\usepackage[T1]{fontenc}
\usepackage[margin=0.5in]{geometry}
\usepackage{parskip}
\usepackage{amsmath}

\pagestyle{fancy}
\fancyhead[l]{Andrew O'Dwyer}
\fancyhead[c]{Ritangle Question 13}
\fancyhead[r]{\today}

\begin{document}

\section{Problem Description}
\huge
Between 0 and 100 hours (not including 100:00:00) on 
stopwatch how many times will exactly 2 unique digits 
be displayed.
\section{Problem Solution}

The Digits ranges are like so:
\begin{align*}
    &Digit_{H,0}: {0-9}\\
    &Digit_{H,1}: {0-9}\\
    &Digit_{M,0}: {0-5}\\
    &Digit_{M,1}: {0-9}\\
    &Digit_{S,0}: {0-5}\\
    &Digit_{S,1}: {0-9}\\
\end{align*}
So there are 4 between 0 and 9 AND there are 2 between 0 and 5.

For the combinations where the digits are less than 6 there are
$2^6 \times \binom{6}{2} = 64 \times 15 = 960$ combinations. 
This is because for a unique pair, each digit can be 1 of 2
states, so there are $2^6$ combinations. And there are 15 
unique pairs.

For the combinations where one of the digits is >=6 and the
other is <=6 there are 0
$2^4 \times 4 \times 6 = 384$ combinations because for the
digits between 6 and 9 there are 4 with 4 different numbers and
for the digits between 0 and 5 there are 2 with 6 different
numbers however these 2 must be the same.

For the combination where both of the digits are >=6 then
there are $2^6 \times \binom{4}{2} = 64 \times 6 = 384$
combinations.

So in total there are 960 + 384 + 384 = 1728 combinations
\end{document}
